\documentclass[11pt]{article}

\usepackage[margin=1in]{geometry}
\usepackage[T1]{fontenc}
\usepackage[utf8]{inputenc}
\usepackage{lmodern}
\usepackage{microtype}
\usepackage{amsmath,amssymb,amsthm}
\usepackage{mathtools}
\usepackage{xcolor}
\usepackage{hyperref}
\usepackage{booktabs}
\usepackage{enumitem}
\usepackage{listings}

\hypersetup{
  colorlinks=true,
  linkcolor=blue!50!black,
  citecolor=blue!50!black,
  urlcolor=blue!50!black
}

\lstdefinelanguage{Lean4}{
  morekeywords={
    import,namespace,end,open,section,noncomputable,abbrev,structure,inductive,
    def,theorem,lemma,where,match,with,if,then,else,by,fun,let,in,return,do,
    simp,calc,have,show,exact,rcases,cases,constructor,intro,apply,
    deriving,instance,class,extends,private,protected,mutual,partial,
    Type,Prop,Sort,Nat,Bool,true,false,sorry,axiom
  },
  sensitive=true,
  morecomment=[l]{--},
  morecomment=[s]{/-}{-/},
  morestring=[b]"
}

\lstset{
  language=Lean4,
  basicstyle=\ttfamily\small,
  keywordstyle=\bfseries\color{blue!50!black},
  commentstyle=\itshape\color{green!40!black},
  stringstyle=\color{red!50!black},
  showstringspaces=false,
  breaklines=true,
  frame=single,
  framerule=0.3pt,
  xleftmargin=0.5em,
  xrightmargin=0.5em,
  keepspaces=true,
  columns=fullflexible,
  escapeinside={@}{@}
}

% Theorem environments
\newtheorem{theorem}{Theorem}[section]
\newtheorem{lemma}[theorem]{Lemma}
\newtheorem{proposition}[theorem]{Proposition}
\newtheorem{corollary}[theorem]{Corollary}
\theoremstyle{definition}
\newtheorem{definition}[theorem]{Definition}
\newtheorem{example}[theorem]{Example}
\theoremstyle{remark}
\newtheorem{remark}[theorem]{Remark}

% Convenience
\newcommand{\MORK}{\textup{MORK}}
\newcommand{\MeTTa}{\textup{MeTTa}}
\newcommand{\MeTTaIL}{\textup{MeTTaIL}}
\newcommand{\PathMap}{\textup{PathMap}}

\title{MORK, MM2, and PathMap: A Lean~4 Formalization\\
\large Syntax, Space Semantics, Trie Data Structure,\\
Conformance Testing, and Cross-System Bridges}
\author{Zar \and Oru\v{z}i\footnote{%
``Oru\v{z}i'' denotes AI-assisted collaboration
using Claude Code (Anthropic), ChatGPT, and Codex (OpenAI).}}
\date{March 2026}

\begin{document}
\maketitle

\begin{abstract}
We present a Lean~4 formalization of the \MORK{} execution engine, its
MM2 (Minimal \MeTTa{}~2) language, and the \PathMap{} trie data structure
that underpins it.
The formalization covers MM2 syntax, space semantics with pattern matching
(including recursive expression matching), a three-phase priority-scheduled
execution protocol, the \PathMap{} algebraic lattice hierarchy with a
concrete finite trie instance (\texttt{FTrie\,V}) satisfying
\texttt{PathMapQuantale}, a bridge to \MeTTaIL{}'s declarative reduction
semantics, and a bridge to MQ-calculus COMM non-determinism.
A conformance test suite of 47 kernel-checked theorems is verified against
the \MORK{} CLI (\texttt{mork run}) on 9+ \texttt{.mm2} test programs.
The \MORK{}/MM2 development comprises 9{,}456 lines across 17 files with
0~sorries and 0~custom axioms.
Since the initial formalization, the development has grown from 8 files
(1{,}911 lines) to 17 files (9{,}456 lines), adding:
a faithful work-queue scheduler with read-copy semantics,
fold-level aggregation (\texttt{FoldAggregator}: selectAll, selectFirst,
count, sum), an idempotent \texttt{head} sink, sink persistence theorems,
N-ary fold non-determinism generalization, an execution boundary
packaging (\texttt{morkTranslatable}), full scheduler-level
correspondence between computable and spec-level work-queue execution
(including multi-match firing, bounded-run equivalence, and an explicit
\texttt{WorkQueueInvariant} bundle), source-side input semantics
covering \texttt{BTM}, \texttt{==} (equality constraint), and
\texttt{!=} (inequality/exclusion constraint) source factors,
a multi-premise \MeTTaIL{} bridge with \texttt{PremiseChain} induction
and freshness guards (\texttt{SourceGuard}), concrete PLN inference
rules (transitivity, modus ponens, abduction, guarded deduction) as
end-to-end GSLT routing witnesses, and a collection congruence bridge
(\texttt{congElem}) proving that sub-element rewriting inside
collections factors through a Huet zipper/lens decomposition
(\texttt{AtomZipper}, \texttt{LensRel}) and corresponds to \MORK{}
rule firings via ground atom self-matching,
proven \texttt{matchAtom} round-trip theorems for compound command
patterns (closing the \texttt{add-atom}/\texttt{remove-atom} MM2
compiler invariant gap),
a pattern normalization function (\texttt{normalize}) that eliminates
\texttt{.subst} nodes via \texttt{openBVar},
an extended \MeTTaIL{} bridge (\texttt{ExtendedBridge.lean})
proving that rewrite rules whose RHS contains \texttt{.subst} nodes
or rest-variable collections still fire through \MORK{}, by showing
that \texttt{applyBindings} always produces \texttt{morkTranslatable}
output (\texttt{morkTranslatable\_applyBindings}),
and a runtime-kernel packaging layer (\texttt{RuntimeKernel.lean},
\texttt{RuntimeResource.lean}) that bridges the \MORK{}/MM2 execution
triad to the \texttt{EffectClass} hierarchy from \texttt{mettail-core}
and classifies current proved fragments by effect class
(\texttt{pureStructural}, \texttt{readOnlyLookup}, \texttt{writesState}),
resource class (\texttt{defaultAtomSpace}), and backend provenance,
with honest descriptors for not-yet-formalized resources (named
atomspaces, Z3 solver, ACT external files).
The \PathMap{} library adds 1{,}262 lines across 5 core files with
54~theorems/instances.
All proofs are checked by Lean~4.28 with Mathlib~v4.28.
\end{abstract}

\tableofcontents

%% ============================================================
\section{Introduction}
\label{sec:intro}
%% ============================================================

\MORK{} (``\MeTTa{} Optimal Reduction Kernel'') is a parallel metagraph
rewriting engine that executes programs written in MM2, a forward-chaining
term-rewriting language over spaces of atoms~\cite{MeredithStayOSLF2020}.
\MORK{} uses \PathMap{} as its underlying data structure for efficient
graph and path representation, and serves as the execution substrate for
\MeTTa{} implementations.

This note documents a Lean~4 formalization of the \MORK{}/MM2/\PathMap{} stack,
covering:
\begin{enumerate}[nosep]
  \item MM2 syntax: exec rules with prioritized pattern-template pairs,
        including \texttt{head} (idempotent add) sinks and fold-level
        aggregators (\S\ref{sec:syntax});
  \item Space semantics: finite sets of ground atoms with substitution-based
        pattern matching and sink application (\S\ref{sec:space});
  \item A relational match specification with soundness and completeness
        (\S\ref{sec:matchspec});
  \item A three-phase execution protocol (unfold/base/fold) with
        priority-band ordering and fold-level aggregation
        (\S\ref{sec:threephase});
  \item A faithful work-queue scheduler with read-copy semantics, exec-fact
        extraction, and ground self-respawn (\S\ref{sec:workqueue});
  \item Sink persistence: atoms persist through sink pipelines unless
        explicitly removed (\S\ref{sec:persistence});
  \item The \PathMap{} algebraic hierarchy: lattice, distributive lattice,
        and quantale, with a concrete finite trie instance (\S\ref{sec:pathmap});
  \item Algebraic correspondence between \MORK{} space transitions and
        \PathMap{} join/subtract (\S\ref{sec:pathmapbridge});
  \item A \MeTTaIL{} bridge: zero-premise, multi-premise,
        freshness-guarded, and collection congruence (\texttt{congElem})
        declarative reductions imply \MORK{} rule firings---the
        \texttt{congElem} bridge uses a Huet zipper/lens factorization
        (\texttt{LensRel}) as the primary congruence surface---with
        concrete PLN inference rules as end-to-end
        routing witnesses (\S\ref{sec:mettail});
  \item An MQ-calculus bridge: binary and N-ary fold non-determinism
        corresponds to \texttt{CommReduction} (\S\ref{sec:mqbridge});
  \item A conformance test suite verified against the \MORK{} CLI
        (\S\ref{sec:conformance});
  \item An execution boundary packaging (\texttt{morkTranslatable}) that
        delineates the proven fragment (\S\ref{sec:execboundary}).
\end{enumerate}

\noindent
Tables~\ref{tab:morkfiles} and~\ref{tab:pathmapfiles} summarize the
formal development.

\begin{table}[ht]
\centering
\caption{\MORK{}/MM2 formalization files.}
\label{tab:morkfiles}
\begin{tabular}{@{}llrl@{}}
\toprule
\textbf{File} & \textbf{Key Content} & \textbf{Lines} \\
\midrule
\texttt{Syntax.lean}              & \texttt{Sink}, \texttt{FoldAggregator}, \texttt{ExecRule}, \texttt{SourceFactor}, \texttt{InputSpec} & 264 \\
\texttt{Space.lean}               & \texttt{Space}, \texttt{matchAtom}, \texttt{applySink}, source-aware matching & 432 \\
\texttt{MatchSpec.lean}           & \texttt{MatchAtomRel}, soundness/completeness & 133 \\
\texttt{ThreePhaseExec.lean}      & Phase protocol, aggregator theorems, canaries & 381 \\
\texttt{ThreePhaseRefinement.lean}& Sink persistence, phase$\leftrightarrow$scheduler & 230 \\
\texttt{WorkQueueOrder.lean}      & \texttt{atomKey}, shortlex ordering, \texttt{lexLt\_iff\_lex} & 241 \\
\texttt{WorkQueueExec.lean}       & Read-copy scheduler, source-aware firing, correspondence & 1{,}192 \\
\texttt{MORKCommBridge.lean}      & MQ correspondence, N-ary generalization & 300 \\
\texttt{PathMapBridge.lean}       & Base/fold = \texttt{pjoin}/\texttt{psubtract} & 174 \\
\texttt{MeTTaILBridge.lean}       & Multi-premise bridge, \texttt{PremiseChain}, ext bridge & 1{,}624 \\
\texttt{PLNRules.lean}            & Concrete PLN rules, guarded bridge instantiation & 383 \\
\texttt{Conformance.lean}         & Computable evaluator, source conformance, correspondence & 1{,}694 \\
\texttt{AtomZipper.lean}          & Huet zipper, \texttt{LensRel}, lens laws, collection specialization & 283 \\
\texttt{CollectionBridge.lean}    & Zipper/lens \texttt{congElem} bridge, \texttt{declReduces\_implies\_mork\_fire\_full} & 319 \\
\texttt{ExecutionBoundary.lean}   & \texttt{morkTranslatable} + scheduler + source + ext + zipper/lens re-exports & 439 \\
\midrule
\textbf{Total} & & \textbf{8{,}089} \\
\bottomrule
\end{tabular}
\end{table}

\begin{table}[ht]
\centering
\caption{\PathMap{} library files (core theory).}
\label{tab:pathmapfiles}
\begin{tabular}{@{}llrl@{}}
\toprule
\textbf{File} & \textbf{Key Content} & \textbf{Lines} & \textbf{Thms/Inst} \\
\midrule
\texttt{Core.lean}              & \texttt{AlgebraicResult}, lattice hierarchy, \texttt{Finset} instance & 440 & 5 / 13 \\
\texttt{CoinductiveTrie.lean}   & \texttt{CTrie V}, pointwise operations, bisimilarity laws & 226 & 30 \\
\texttt{FiniteTrie.lean}        & \texttt{FTrie V}, computable trie operations & 224 & 5 \\
\texttt{TrieRefinement.lean}    & \texttt{Sorted}, \texttt{join\_lookup}, \texttt{toCTrie} & 243 & 5 \\
\texttt{TriePathMapInstance.lean} & \texttt{PathMapQuantale (FTrie V)}, \texttt{JoinIdem}, \texttt{MeetIdem} & 129 & 0 / 4 \\
\midrule
\textbf{Total} & & \textbf{1{,}262} & \textbf{45 + 17 inst} \\
\bottomrule
\end{tabular}
\end{table}

All \MORK{}/MM2 files reside under
\texttt{Mettapedia/Languages/ProcessCalculi/MORK/}.
The \PathMap{} library is at \texttt{Mettapedia/OSLF/PathMap/}.

%% ============================================================
\section{MM2 Syntax}
\label{sec:syntax}
%% ============================================================

MM2 programs consist of \emph{exec rules} that match atoms in a space and
produce add/remove mutations.
The grammar, as formalized in Lean following the MORK wiki spec, is:
\begin{center}
\begin{tabular}{@{}rcl@{}}
\textit{ExecRule}  & $::=$ & \texttt{(exec (priority $n$ name) Pattern Template)} \\
\textit{Pattern}   & $::=$ & \texttt{(, Atom$_1$ \ldots{} Atom$_n$)} \\
\textit{Template}  & $::=$ & \texttt{(O Sink$_1$ \ldots{} Sink$_n$)} \\
\textit{Sink}      & $::=$ & \texttt{(+ Atom)} $\mid$ \texttt{(- Atom)} $\mid$ \texttt{(head Atom)} \\
\end{tabular}
\end{center}

\noindent
A \emph{pattern} is a conjunction: all atoms must be present in the space
for the rule to fire.
A \emph{template} is a list of sinks applied simultaneously: \texttt{(+ a)}
inserts atom~\texttt{a}, \texttt{(- a)} removes it, and \texttt{(head a)}
performs an idempotent add (first-wins: no-op if already present).

The Lean formalization uses the shared \texttt{Atom} type from
\texttt{MeTTa.Core.Atom}, so \MORK{} S-expressions are a subset of \MeTTa{}
atoms:

\begin{lstlisting}[caption={MM2 syntax types (\texttt{Syntax.lean}).}]
structure Pattern where
  atoms : List Atom
  deriving Repr, DecidableEq

inductive Sink where
  | add    : Atom -> Sink   -- (+ a)
  | remove : Atom -> Sink   -- (- a)
  | head   : Atom -> Sink   -- (head a): idempotent add
  deriving Repr, DecidableEq

structure Template where
  sinks : List Sink
  deriving Repr, DecidableEq

structure ExecRule where
  priority : Nat
  name     : String
  pat      : Pattern
  tmpl     : Template
  deriving Repr, DecidableEq
\end{lstlisting}

\subsection{Fold-Level Aggregation}

Sub-results assembled during the fold phase are governed by a
\texttt{FoldAggregator}:

\begin{lstlisting}[caption={Fold aggregator (\texttt{Syntax.lean}).}]
inductive FoldAggregator where
  | selectAll   : FoldAggregator  -- non-deterministic (default)
  | selectFirst : FoldAggregator  -- head only
  | count       : FoldAggregator  -- (.grounded (.int N))
  | sum         : FoldAggregator  -- sum of integer sub-results
  deriving Repr, DecidableEq

def applyAggregator (agg : FoldAggregator)
    (subResults : List Atom) : Option Atom :=
  match agg with
  | .selectAll   => subResults.head?
  | .selectFirst => subResults.head?
  | .count       =>
      some (.grounded (.int subResults.length))
  | .sum         =>
      let vals := subResults.filterMap extractInt
      some (.grounded (.int (vals.foldl (. + .) 0)))
\end{lstlisting}

\noindent
The \texttt{selectAll} and \texttt{selectFirst} cases both return
\texttt{head?} at this level; the non-determinism of \texttt{selectAll} is
captured by \texttt{NaryFoldPicksSubResult} in \texttt{MORKCommBridge.lean},
not by \texttt{applyAggregator}.

\begin{theorem}[\texttt{applyAggregator\_count}]
\texttt{applyAggregator .count rs = some (.grounded (.int rs.length))}.
\end{theorem}

\begin{theorem}[\texttt{applyAggregator\_selectFirst}]
For $\mathit{rs} = a :: \mathit{tl}$:
\texttt{applyAggregator .selectFirst rs = some a}.
\end{theorem}

\begin{theorem}[\texttt{applyAggregator\_sum\_nil}]
\texttt{applyAggregator .sum [] = some (.grounded (.int 0))}.
\end{theorem}

%% ============================================================
\section{Space Semantics and Matching}
\label{sec:space}
%% ============================================================

\subsection{Space and Substitution}

A \MORK{} \emph{space} is a finite set of ground atoms.
The Rust implementation uses \texttt{PathMap<()>}; the formalization abstracts
this to \texttt{Finset Atom} for tractable reasoning:

\begin{lstlisting}[caption={Space and substitution (\texttt{Space.lean}).}]
abbrev Space := Finset Atom
abbrev Subst := List (String * Atom)
\end{lstlisting}

Substitutions are association lists with first-match-wins lookup via
\texttt{Subst.lookup}.

\subsection{Pattern Matching}

The spec-level function \texttt{matchAtom} implements first-order
(non-unification) pattern matching on three atom kinds:

\begin{itemize}[nosep]
  \item \textbf{Variable} (\texttt{.var v}): if \texttt{v} is already bound
    in the substitution, the concrete atom must equal the bound value;
    otherwise, \texttt{v} is freshly bound.
  \item \textbf{Symbol} (\texttt{.symbol s}): must match an identical symbol.
  \item \textbf{Grounded} (\texttt{.grounded g}): must match an identical
    grounded value.
  \item \textbf{Expression} (\texttt{.expression es}): returns \texttt{none}
    in the spec-level \texttt{matchAtom} (see \S\ref{sec:conformance} for
    the computable evaluator that handles expressions).
\end{itemize}

\noindent
The conjunctive match \texttt{matchPattern} threads the substitution through
a list of pattern atoms, matching each against the space while tracking
consumed atoms to prevent double-matching.

\subsection{Sink Application and Rule Firing}

A single sink mutates the space:

\begin{lstlisting}[caption={Sink application (\texttt{Space.lean}).}]
def applySink (s : Space) (sig : Subst) (sink : Sink)
    : Space :=
  match sink with
  | .add a    =>
    let a' := applySubst sig a
    if isGroundAtom a' then s @$\cup$@ {a'} else s
  | .remove a => s.erase (applySubst sig a)
  | .head a   =>
    let a' := applySubst sig a
    if isGroundAtom a' then s @$\cup$@ {a'} else s
\end{lstlisting}

\noindent
Non-ground add results are silently dropped (a guard for soundness).
In the \texttt{Finset} model, \texttt{head} is definitionally equal to
\texttt{add} (union is idempotent).  The distinction is operative in the
computable reference evaluator (\texttt{Conformance.lean}), which uses
\texttt{List Atom} where idempotence must be checked explicitly via
\texttt{contains}.
A template's sinks are applied left-to-right via \texttt{applySinks},
and \texttt{fireRule} composes matching with template application.

\subsection{Structural Lemmas}

Four theorems characterize sink effects:

\begin{theorem}[\texttt{applySink\_add\_card}]
Adding a ground atom not already in the space increases cardinality by~1.
\end{theorem}

\begin{theorem}[\texttt{applySink\_remove\_card}]
Removing an atom present in the space decreases cardinality by~1.
\end{theorem}

\begin{theorem}[\texttt{applySinks\_empty}]
An empty template is the identity on the space.
\end{theorem}

\begin{theorem}[\texttt{applySink\_add\_mem}]
Adding an atom preserves membership of all existing atoms.
\end{theorem}

%% ============================================================
\section{Match Specification}
\label{sec:matchspec}
%% ============================================================

The file \texttt{MatchSpec.lean} introduces a relational specification for
pattern matching, separating the ``what'' from the ``how'':

\begin{lstlisting}[caption={Match relation (\texttt{MatchSpec.lean}).}]
inductive MatchAtomRel : Subst -> Atom -> Atom
    -> Subst -> Prop where
  | var_fresh  (hlookup : sig.lookup v = none) :
      MatchAtomRel sig (.var v) a ((v, a) :: sig)
  | var_bound  (hlookup : sig.lookup v = some a) :
      MatchAtomRel sig (.var v) a sig
  | symbol :  MatchAtomRel sig (.symbol s) (.symbol s) sig
  | grounded : MatchAtomRel sig (.grounded g) (.grounded g) sig
\end{lstlisting}

The key result is an exact characterization:

\begin{theorem}[\texttt{matchAtom\_iff}]
For all substitutions $\sigma$, pattern atoms $p$, concrete atoms $c$, and
result substitutions $\sigma'$:
$
  \mathtt{matchAtom}\;\sigma\;p\;c = \mathtt{some}\;\sigma'
  \;\iff\;
  \mathtt{MatchAtomRel}\;\sigma\;p\;c\;\sigma'
$
\end{theorem}

\noindent
This is proved via separate soundness and completeness theorems.
Expression patterns are excluded: \texttt{matchAtom\_expression\_pattern\_none}
proves that \texttt{matchAtom} on \texttt{.expression} always returns
\texttt{none}.

%% ============================================================
\section{Three-Phase Execution}
\label{sec:threephase}
%% ============================================================

\MORK{}'s execution protocol schedules rules by priority into three
non-overlapping bands:

\begin{definition}[Phase Bands]
\begin{center}
\begin{tabular}{@{}lll@{}}
\toprule
\textbf{Phase} & \textbf{Priority Range} & \textbf{Role} \\
\midrule
Unfold & $0$--$31$   & Spawn $N$ sub-queries + wait token \\
Base   & $32$--$63$  & Resolve leaf queries directly \\
Fold   & $64$--$95$  & Wait for sub-results, assemble final result \\
\bottomrule
\end{tabular}
\end{center}
\end{definition}

Each phase has a dedicated step type (\texttt{UnfoldStep}, \texttt{BaseStep},
\texttt{FoldStep}) with a proof field witnessing that the step's priority
falls within the correct band.
The \texttt{FoldStep} structure carries an \texttt{aggregator :
FoldAggregator} field (default \texttt{.selectAll}) governing how
sub-results are assembled.
Space transitions \texttt{applyUnfold}, \texttt{applyBase}, and
\texttt{applyFold} implement each phase.

\begin{theorem}[\texttt{phase\_ranges\_disjoint}]
No priority belongs to two distinct phases.
\end{theorem}

\begin{theorem}[\texttt{phase\_priority\_monotone}]
For any priorities $n_u$, $n_b$, $n_f$ in the unfold, base, and fold
phases respectively: $n_u < n_b < n_f$.
\end{theorem}

\begin{theorem}[\texttt{foldStep\_default\_aggregator}]
For a fold step with default aggregator (\texttt{.selectAll}):
\texttt{applyAggregator fold.aggregator fold.subResults =
fold.subResults.head?}.
\end{theorem}

%% ============================================================
\section{Work-Queue Scheduler}
\label{sec:workqueue}
%% ============================================================

The work-queue scheduler (\texttt{WorkQueueExec.lean}, 1{,}143~lines)
formalizes \MORK{}'s actual runtime model from \texttt{mork\_backend.rs}.
This replaces the abstract three-phase view with the concrete execution
loop:

\begin{enumerate}[nosep]
\item \textbf{Pop} the next \texttt{(exec ...)} fact from the space
  (sorted by \texttt{atomKey}: serialized key with shortlex ordering).
\item \textbf{Remove} it from live space.
\item \textbf{Read copy}: $\mathrm{live\_after\_remove} \cup
  \{\mathrm{exec\_fact}\}$---re-insert so the rule can match itself.
\item \textbf{Match} the pattern against the read copy.
\item \textbf{Apply} template sinks to live space (not to read copy).
\item \textbf{Repeat} until no exec facts remain or step limit reached.
\end{enumerate}

\noindent
The read-copy semantics are critical: step~3 ensures that a rule can
match its own exec fact (required for self-referential patterns).

\begin{theorem}[\texttt{readCopy\_mem\_exec}]
The exec fact is always a member of its own read copy.
\end{theorem}

\begin{theorem}[\texttt{readCopy\_eq\_of\_mem}]
Remove + re-insert = identity on membership: if $a \in s$ then
$a \in (s \setminus \{a\}) \cup \{a\}$.
\end{theorem}

\subsection{Exec-Fact Extraction}

The function \texttt{extractExecFact} parses an \texttt{(exec loc pat tpl)}
atom into a structured \texttt{ExecFact} with priority, name, pattern, and
template fields.  Sink parsing handles \texttt{(+ ...)}, \texttt{(- ...)},
and \texttt{(head ...)} forms.

\subsection{Ordering}

\texttt{WorkQueueOrder.lean} (119~lines) defines the
\texttt{atomKey} serialization and \texttt{lexLt} comparison used to
determine exec-fact firing order:

\begin{theorem}[\texttt{order\_p0\_lt\_p1}]
\texttt{atomKey} orders priority~0 before priority~1.
\end{theorem}

\begin{theorem}[\texttt{order\_1\_lt\_10}]
Shortlex ordering: single-digit priorities fire before double-digit.
\end{theorem}

\subsection{Ground Self-Respawn}

A canary theorem exercises the scheduler with a ground rule that
re-adds its own exec fact via a \texttt{(+ ...)} sink:

\begin{theorem}[\texttt{canary8\_ground\_self\_respawn}]
A ground self-respawn rule fires exactly once---the second step finds
no new match since the exec fact is already present in the space.
\end{theorem}

\begin{theorem}[\texttt{canary8\_second\_step\_no\_match}]
Running the scheduler for 2 steps on the self-respawn rule yields
the fired result plus the original exec fact, with step count~2.
\end{theorem}

\subsection{Scheduler-Level Correspondence}
\label{sec:schedulercorr}

The computable work-queue (\texttt{cWorkQueueStep}, \texttt{cWorkQueueRunN})
operates on \texttt{List Atom}, while the spec-level definitions
(\texttt{workQueueStep}, \texttt{workQueueRunN}) operate on
\texttt{Finset Atom}.  The scheduler correspondence chain establishes
that these agree at the \texttt{toFinset} level, building from per-piece
soundness through full bounded-run equivalence.

\paragraph{Per-sink correspondence.}
Each computable sink step (\texttt{capplySinkStep}) agrees with the
spec-level \texttt{applySink}, provided a \texttt{NodupSafe} invariant
holds (no duplicates at each remove step during the foldl over sinks).
This is the base layer.

\paragraph{Multi-match foldl correspondence.}
The general case lifts per-sink correspondence through the outer foldl
over match results.  The predicate \texttt{FoldNodupSafe} ensures
\texttt{NodupSafe} at every step of this outer fold:

\begin{definition}[\texttt{FoldNodupSafe}]
Given an accumulator $\mathit{acc}$, a list of match results
$\mathit{ms}$, and a template $\mathit{tmpl}$:
$\texttt{FoldNodupSafe}(\mathit{acc}, \mathit{ms}, \mathit{tmpl})$
holds iff for every $i < |\mathit{ms}|$, the accumulator after
folding the first $i$ results satisfies \texttt{NodupSafe} for
the $i$-th substitution and the template's sinks.
\end{definition}

\begin{theorem}[\texttt{foldl\_capplySinks\_toFinset}]
If the computable and spec match results have the same substitutions
in the same order and \texttt{FoldNodupSafe} holds, then the foldls
produce corresponding results at the \texttt{Finset} level.
\end{theorem}

\noindent
The alignment hypothesis (same substitutions in same order) is explicit
because \texttt{Finset.toList} has unspecified order.  For patterns that
match at most once (the common MORK case), this is trivially satisfied.

\paragraph{General firing correspondence.}

\begin{theorem}[\texttt{cFireExecFact\_toFinset}]
For any exec fact $\mathit{ef}$ with $\mathit{ef}.\mathit{atom} \in s$,
if match results align and \texttt{FoldNodupSafe} holds,
then $\texttt{cFireExecFact}(s, \mathit{ef}).\texttt{toFinset} =
\texttt{fireExecFact}(s.\texttt{toFinset}, \mathit{ef})$.
\end{theorem}

\paragraph{Step and bounded-run correspondence.}

\begin{theorem}[\texttt{cWorkQueueStep\_toFinset}]
If the \texttt{WorkQueueInvariant} holds (bundling \texttt{Nodup},
\texttt{KeyInjective}, and per-firing alignment), then the computable
and spec work-queue steps agree:
$(\texttt{cWorkQueueStep}\;s).\texttt{map}\;\texttt{toFinset}
  = \texttt{workQueueStep}\;(s.\texttt{toFinset})$.
\end{theorem}

\begin{theorem}[\texttt{cWorkQueueRunN\_toFinset}]
If the \texttt{WorkQueueInvariant} holds at every computably reachable
state (formalized via the \texttt{CReachable} inductive), then
$n$-step bounded execution agrees:
\[
  (\texttt{cWorkQueueRunN}\;n\;s).1.\texttt{toFinset}
  = (\texttt{workQueueRunN}\;n\;(s.\texttt{toFinset})).1
  \quad\text{and}\quad
  (\texttt{cWorkQueueRunN}\;n\;s).2
  = (\texttt{workQueueRunN}\;n\;(s.\texttt{toFinset})).2
\]
\end{theorem}

\noindent
The invariant-maintenance burden is explicit: for finite-state MORK
programs (the common case), it can be verified by exhaustive enumeration.

\subsection{Source-Side Input Semantics}
\label{sec:source}

MORK exec rules support two input modes: \emph{compat mode}
\texttt{(, pat$_1$ \ldots pat$_n$)}, where all pattern atoms are matched
against the workspace, and \emph{explicit source mode}
\texttt{(I src$_1$ \ldots src$_n$)}, where each source factor specifies
how its pattern is matched.

The formalization covers three source factor types
(\texttt{SourceFactor} in \texttt{Syntax.lean}):
\begin{itemize}[nosep]
\item \texttt{btm}: Match pattern against the workspace (equivalent to
  compat-mode per-atom matching).
\item \texttt{eqConstraint} (\texttt{==}): Substitute the first
  sub-expression using current bindings, check if the result exists in
  the space, and if so bind the witness to the found atom.
\item \texttt{neqConstraint} (\texttt{!=}): Substitute the first
  sub-expression, \emph{remove} it from the candidate set, then match
  the witness against remaining atoms.  This implements the Rust
  \texttt{CmpSource} with \texttt{cmp=1}.
\end{itemize}

\begin{theorem}[\texttt{fireSourceRule\_compat}]
Source-aware firing on a compat-mode rule agrees with the standard
\texttt{fireRule}: $\texttt{fireSourceRule}\;s\;(r.\texttt{toSourceRule})
= \texttt{fireRule}\;s\;r$.
\end{theorem}

The computable evaluator (\texttt{cmatchInputSpec},
\texttt{cfireSourceRule}, \texttt{cFireSourceExecFact}) mirrors the
spec-level definitions, and 11 kernel-checked conformance tests
(tests S1--S8 in \texttt{Conformance.lean} plus canaries 9--11 in
\texttt{WorkQueueExec.lean}) verify the implementation against
\texttt{mork run} ground truth, covering BTM matching, equality
constraint success/failure, shared variable propagation,
inequality exclusion, and multi-match non-determinism.

%% ============================================================
\section{Sink Persistence}
\label{sec:persistence}
%% ============================================================

\texttt{ThreePhaseRefinement.lean} (212~lines) proves that atoms persist
through sink pipelines unless explicitly removed:

\begin{theorem}[\texttt{applySink\_mem\_of\_mem}]
If $b \in s$ and the sink is not a remove of~$b$ (after substitution),
then $b$ remains after applying the sink.  Handles all three sink
constructors: \texttt{add}, \texttt{remove}, and \texttt{head}.
\end{theorem}

\begin{theorem}[\texttt{applySinks\_mem\_of\_mem}]
If $b \in s$ and no sink in the template removes~$b$, then $b \in
\texttt{applySinks}\;s\;\sigma\;\mathit{tmpl}$.
\end{theorem}

\noindent
The proof proceeds inductively via a helper
\texttt{foldl\_applySink\_mem} that threads the non-removal invariant
through the \texttt{List.foldl} over sinks.

\begin{theorem}[\texttt{applySink\_add\_ground\_mem}]
Adding a ground atom via \texttt{Sink.add} guarantees the substituted
atom is a member of the resulting space.
\end{theorem}

%% ============================================================
\section{PathMap Algebraic Hierarchy}
\label{sec:pathmap}
%% ============================================================

\PathMap{} is a trie-indexed data structure with algebraic operations that
form a quantale.  The formalization provides both the abstract algebraic
hierarchy and a concrete trie implementation.

\subsection{AlgebraicResult and Typeclass Hierarchy}

Operations on \PathMap{} values return an \texttt{AlgebraicResult\,V}, an
inductive type that encodes structural sharing:

\begin{lstlisting}[caption={\texttt{AlgebraicResult} (\texttt{Core.lean}).}]
inductive AlgebraicResult (V : Type u) where
  | none : AlgebraicResult V
  | identity (isSelf : Bool) (isOther : Bool)
      : AlgebraicResult V
  | element : V -> AlgebraicResult V
\end{lstlisting}

\noindent
The \texttt{identity} variant avoids allocation when the result equals one
of the inputs.
Three typeclasses form the hierarchy:

\begin{itemize}[nosep]
  \item \textbf{\texttt{PathMapLattice\,$\alpha$}}: \texttt{pjoin} (union)
    and \texttt{pmeet} (intersection).
  \item \textbf{\texttt{PathMapDistributiveLattice\,$\alpha$}}: adds
    \texttt{psubtract} (set difference).
  \item \textbf{\texttt{PathMapQuantale\,$\alpha$}}: adds \texttt{prestrict}
    (prefix filter).
\end{itemize}

\noindent
Algebraic law classes (\texttt{JoinComm}, \texttt{JoinIdem},
\texttt{MeetComm}, \texttt{MeetIdem}, \texttt{Absorption}) are defined at
the value level.

\subsection{Finset Instance}

\texttt{Core.lean} provides a \texttt{PathMapQuantale (Finset\,$\alpha$)}
instance with all algebraic laws proved:

\begin{itemize}[nosep]
  \item \texttt{pjoin} = union with identity detection (\texttt{a = b},
    \texttt{a $\subseteq$ b}, \texttt{b $\subseteq$ a}).
  \item \texttt{pmeet} = intersection; returns \texttt{.none} if empty.
  \item \texttt{psubtract} = set difference; returns \texttt{.none} if empty.
  \item \texttt{prestrict} = subset check then intersection.
\end{itemize}

\noindent
This is the instance used by \texttt{PathMapBridge.lean} for
\texttt{Space = Finset Atom}.

\subsection{Coinductive Trie (CTrie)}

The file \texttt{CoinductiveTrie.lean} defines the semantic reference model:

\begin{lstlisting}[caption={Coinductive trie (\texttt{CoinductiveTrie.lean}).}]
abbrev CTrie V := List UInt8 -> Option V
\end{lstlisting}

\noindent
A trie IS its lookup function.
Pointwise operations (\texttt{union}, \texttt{inter}, \texttt{diff},
\texttt{restrict}) are defined by lifting \texttt{Option} combinators, and
30 theorems prove bisimilarity laws (idempotence, commutativity, identity
elements, associativity).

\subsection{Finite Trie (FTrie)}

The concrete, computable trie:

\begin{lstlisting}[caption={Finite trie (\texttt{FiniteTrie.lean}).}]
inductive FTrie (V : Type u) where
  | empty : FTrie V
  | node (val : Option V)
         (children : List (UInt8 * FTrie V))
      : FTrie V
\end{lstlisting}

\noindent
Operations (\texttt{lookup}, \texttt{join}, \texttt{meet}, \texttt{subtract},
\texttt{restrict}) use mutual recursion with child-list helpers.
A \texttt{Sorted} predicate (\texttt{TrieRefinement.lean}) ensures child
keys are ordered, and refinement theorems prove that \texttt{FTrie}
operations commute with \texttt{CTrie} operations:

\begin{theorem}[\texttt{join\_lookup}]
$(t_1.\mathtt{join}\;t_2).\mathtt{lookup}\;p
  = t_1.\mathtt{lookup}\;p \mathbin{<\!|>} t_2.\mathtt{lookup}\;p$
\quad (for \texttt{Sorted} tries).
\end{theorem}

\subsection{PathMapQuantale Instance for FTrie}

\texttt{TriePathMapInstance.lean} provides:

\begin{lstlisting}[caption={\texttt{PathMapQuantale} instance for \texttt{FTrie}.}]
instance : PathMapQuantale (FTrie V) where
  pjoin a b :=
    if beqFTrie a b then .identity true true
    else .element (a.join b)
  pmeet a b := ...   -- similar with emptiness check
  psubtract a b := ...
  prestrict a b := ...
\end{lstlisting}

\noindent
\texttt{JoinIdem} and \texttt{MeetIdem} are proved (via
\texttt{beqFTrie\_refl}).
Full commutativity holds for \texttt{V = Unit} (the \MORK{} use case where
\texttt{PathMap<()>} stores set membership).

%% ============================================================
\section{MORK--PathMap Bridge}
\label{sec:pathmapbridge}
%% ============================================================

\texttt{Space = Finset Atom} inherits the \texttt{PathMapQuantale} instance.
The bridge file proves that \MORK{}'s space transitions decompose into
algebraic operations:

\begin{theorem}[\texttt{pjoin\_singleton\_resolves}]
$\mathtt{PJ}\;s\;\{a\}$ resolves to $s \cup \{a\}$.
\end{theorem}

\begin{theorem}[\texttt{psubtract\_singleton\_erase}]
$\mathtt{PS}\;s\;\{a\}$ resolves (with $\emptyset$ as default) to
$s \setminus \{a\}$.
\end{theorem}

\begin{theorem}[\texttt{applyBase\_eq\_lattice\_ops}]
\texttt{applyBase} = \texttt{PS} (remove query) then \texttt{PJ} (add result).
\end{theorem}

\begin{theorem}[\texttt{applyFold\_eq\_lattice\_ops}]
\texttt{applyFold} = chained \texttt{PS} (remove wait + sub-results) then
\texttt{PJ} (add assembled result).
\end{theorem}

These theorems confirm that \MORK{}'s operational transitions are expressible
as algebraic operations on the \PathMap{} lattice, connecting the execution
engine to the data-structural semantics.

%% ============================================================
\section{MeTTaIL Bridge}
\label{sec:mettail}
%% ============================================================

The \MeTTaIL{} bridge (\texttt{MeTTaILBridge.lean}, 1{,}635~lines, plus
\texttt{CollectionBridge.lean}, 217~lines) connects \MORK{}'s execution
model to the declarative reduction semantics of \MeTTaIL{}
(\texttt{DeclReduces} and \texttt{DeclReducesWithPremises}).

The mutual recursive function \texttt{morkPatternToAtom} translates
\MeTTaIL{} patterns to \MORK{} atoms.
The predicate \texttt{morkTranslatable} identifies soundly translatable
patterns (no beta-redex, no rest-variables).

\subsection{Zero-Premise Bridge}

\begin{theorem}[\texttt{declReduces\_implies\_mork\_fire}]
\label{thm:bridge}
If a \MeTTaIL{} declarative reduction applies a rule (with \texttt{r.left =
.fvar x}, translatable RHS, ground result, and no collection-element
rewriting), then the corresponding \MORK{} exec rule fires and the result
space contains the translated output.
\end{theorem}

\subsection{Multi-Premise Bridge via PremiseChain}

\MeTTaIL{} rules may have \emph{premises}: side-conditions that must be
satisfied for the rule to fire.  The two premise types relevant to \MORK{}
are \texttt{relationQuery} (workspace lookup) and \texttt{freshness}
(variable-distinctness guard).

A \texttt{PremiseChain} is an inductive proof that a list of premises
can be satisfied by workspace atoms:

\begin{lstlisting}[caption={PremiseChain inductive type.}]
inductive PremiseChain (relEnv) (lang) (s : Space)
    : Bindings -> List Premise -> List Atom -> Bindings -> Prop where
  | nil  : PremiseChain relEnv lang s bs [] [] bs
  | cons : (prem : Premise) -> (a : Atom) ->
           a @$\in$@ s ->
           matchAtom (@$\sigma_0$@) (premiseToAtomTotal prem) a = some @$\sigma_1$@ ->
           @$b_1$@ @$\in$@ premiseStepWithEnv relEnv lang @$b_0$@ prem ->
           PremiseChain relEnv lang s @$b_1$@ prems witnesses @$b_f$@ ->
           PremiseChain relEnv lang s @$b_0$@ (prem :: prems)
                                      (a :: witnesses) @$b_f$@
  | guard : (prem : Premise) ->
            premiseToSourceFactor prem = none ->
            @$b_0$@ @$\in$@ premiseStepWithEnv relEnv lang @$b_0$@ prem ->
            PremiseChain relEnv lang s @$b_0$@ prems witnesses @$b_f$@ ->
            PremiseChain relEnv lang s @$b_0$@ (prem :: prems) witnesses @$b_f$@
\end{lstlisting}

\noindent
The \texttt{cons} constructor handles \texttt{relationQuery} premises
(consuming a workspace atom as witness), while \texttt{guard} handles
\texttt{freshness} premises (checking a substitution-level condition without
consuming workspace atoms).

\begin{theorem}[\texttt{declReducesWithPremises\_multi\_implies\_mork\_fireSourceRule}]
\label{thm:multibridge}
If a \texttt{DeclReducesWithPremises} reduction fires using a rule whose
premises are all \texttt{relationQuery} (\texttt{allPremisesTranslatable}),
and a \texttt{PremiseChain} witness exists in the workspace, then the
translated \texttt{SourceExecRule} fires in \texttt{fireSourceRule}.
\end{theorem}

\subsection{Freshness Guards and the Extended Bridge}

\texttt{SourceGuard.freshness} encodes the condition that a variable does not
occur free in a substituted term.  The correspondence theorem bridges
\MeTTaIL{}'s sequential freshness checking to \MORK{}'s guard filtering:

\begin{theorem}[\texttt{freshness\_premise\_correspond}]
If \MeTTaIL{}'s freshness premise succeeds at bindings~$\beta$, then
\texttt{matchSourceGuard ($\texttt{bindingsToSubst}\;\beta$) g = true}
for the translated guard~$g$.
\end{theorem}

\noindent
The extended bridge handles premise lists mixing \texttt{relationQuery}
and \texttt{freshness}:

\begin{theorem}[\texttt{declReducesWithPremises\_ext\_implies\_mork\_fireSourceRule}]
\label{thm:extbridge}
If a \texttt{DeclReducesWithPremises} reduction fires using a rule whose
premises are all ext-translatable (\texttt{allPremisesTranslatableExt}),
a \texttt{PremiseChain} witness exists, and the extracted guards pass at
the final substitution, then the translated \texttt{SourceExecRule} (with
guards) fires in \texttt{fireSourceRule}.
\end{theorem}

\noindent
\textbf{Design note.}
\MeTTaIL{} checks freshness guards \emph{sequentially} at the point each
guard fires in the premise chain, while \MORK{}'s \texttt{fireSourceRule}
checks guards at the \emph{final} substitution.  These semantics differ in
general: extending a substitution can break a previously-satisfied freshness
condition.  The bridge therefore accepts guard satisfaction at the final
substitution as an explicit hypothesis.  For typical rules where guard
variables are fully bound by earlier factor-matching premises (as in all
concrete PLN rules below), this hypothesis is trivially dischargeable.

\subsection{Concrete PLN Rules as GSLT Routing Witnesses}

\texttt{PLNRules.lean} (379~lines, 23~theorems) instantiates the bridge for
concrete PLN inference rules, providing end-to-end GSLT routing witnesses:

\begin{center}
\begin{tabular}{@{}llll@{}}
\toprule
\textbf{Rule} & \textbf{Premises} & \textbf{Conclusion} & \textbf{Guards} \\
\midrule
Transitivity       & Inh(A,B), Inh(B,C)             & Inh(A,C)  & --- \\
Modus ponens       & Impl(P,Q), Eval(P)              & Eval(Q)   & --- \\
Inh$\to$Similarity & Inh(A,B)                        & Sim(A,B)  & --- \\
Abduction          & Inh(A,B), Inh(C,B)              & Sim(A,C)  & --- \\
Guarded trans.     & Inh(A,B), Inh(B,C)              & Inh(A,C)  & A\#C \\
Double-guarded     & Inh(A,B), Inh(B,C)              & Inh(A,C)  & B\#A, B\#C \\
\bottomrule
\end{tabular}
\end{center}

\noindent
All translatability proofs are by \texttt{decide}.
Two language definitions package these rules:
\texttt{plnPremiseLanguageDef} (3~unguarded rules, using the strict bridge)
and \texttt{plnGuardedPremiseLanguageDef} (3~guarded/mixed rules, using the
ext bridge).  Concrete bridge instantiation theorems
(\texttt{pln\_reduces\_implies\_mork\_fireSourceRule} and
\texttt{pln\_guarded\_reduces\_implies\_mork\_fireSourceRule})
demonstrate the full chain from \texttt{DeclReducesWithPremises} to
\texttt{fireSourceRule} for each language definition.

\begin{remark}[Honest Scope]
The multi-premise and extended bridges handle flat rules where the LHS is
a free variable.  Collection-element congruence (\texttt{congElem}) is
handled separately via the collection bridge (\S\ref{sec:collbridge}).
Beta-redex patterns and non-ground results remain outside the formalized
boundary.
\end{remark}

\subsection{Collection Congruence Bridge}
\label{sec:collbridge}

\texttt{CollectionBridge.lean} (319~lines) and \texttt{AtomZipper.lean}
(283~lines) prove that \MeTTaIL{}'s \texttt{congElem} constructor---which
rewrites a single element inside a \texttt{.collection ct elems rest}---corresponds
to a \MORK{}-side space update via a principled zipper/lens factorization.

The bridge is organized in layers:

\paragraph{Layers A--B: Structural infrastructure.}
\texttt{morkPatternToAtomList} commutes with \texttt{List.set} (Layer~A),
and ground atoms match themselves under any substitution (Layer~B):

\begin{lstlisting}
theorem morkPatternToAtomList_set (elems : List ILP)
    (i : Nat) (q' : ILP) :
    morkPatternToAtomList (elems.set i q') =
    (morkPatternToAtomList elems).set i
      (morkPatternToAtom q')

theorem matchAtom_ground_self (@$\sigma$@ : Subst) (a : Atom)
    (hg : isGroundAtom a = true) :
    matchAtom @$\sigma$@ a a = some @$\sigma$@
\end{lstlisting}

\paragraph{Layer C: Ad-hoc collection-replace rule.}
For a \texttt{congElem} step replacing element~$i$ in a collection,
we construct an exec rule that matches the entire old collection
atom and replaces it with the modified one.  Since both atoms are
ground, self-matching and identity substitution guarantee correctness.

\paragraph{Layer D (primary): Huet zipper and lens relation.}
\texttt{AtomZipper.lean} formalizes a concrete Huet zipper over
\texttt{Atom}, following \MORK{}'s Rust \texttt{ExprZipper}:

\begin{lstlisting}
structure AtomCrumb where
  left  : List Atom   -- reversed left siblings
  right : List Atom   -- right siblings

abbrev AtomContext := List AtomCrumb
structure AtomZipper where
  focus : Atom
  ctx   : AtomContext
\end{lstlisting}

\noindent
Navigation (\texttt{descendExprChild?}, \texttt{ascend?}),
path-indexed focus (\texttt{focusAtPath?}), and
path-indexed replacement (\texttt{replaceAtPath?}) are proven
to compose correctly.  A \texttt{LensRel} captures the
get/put/put-put factorization:

\begin{lstlisting}
def LensRel (whole part newPart newWhole : Atom)
    : Prop :=
  @$\exists$@ z : AtomZipper,
    rebuild z = whole @$\wedge$@ z.focus = part @$\wedge$@
    rebuild (replaceFocus z newPart) = newWhole
\end{lstlisting}

\noindent
Three classical lens laws are proven (0~sorries):
\texttt{lensRel\_get\_put} (replace focus with itself is identity),
\texttt{lensRel\_put\_get} (focus of replacement equals new value),
and \texttt{lensRel\_put\_put} (consecutive replacements compose).

The collection specialization \texttt{collection\_lensRel} proves
that replacing element~$i$ in a \texttt{.expression (ctSym :: elems)}
satisfies \texttt{LensRel}, connecting the zipper decomposition to
\MeTTaIL{}'s \texttt{congElem} constructor.

The canonical congruence bridge combines the structural and semantic layers:

\begin{theorem}[\texttt{congElem\_implies\_mork\_zipper\_fire}]
\label{thm:congelemzipper}
For any \texttt{congElem} reduction step where the collection atom is
in the \MORK{} space and both old and new atoms are ground, the
element update factors through \texttt{LensRel} \emph{and} there exists
an exec rule whose firing produces the updated space.
\end{theorem}

\paragraph{Layer E: Semantic-only corollary.}
\texttt{congElem\_implies\_mork\_fire} provides the same exec-rule
firing guarantee without the \texttt{LensRel} structural witness.

\paragraph{Layer F: Full \texttt{DeclReduces} bridge.}
\texttt{declReduces\_implies\_mork\_fire\_full} handles both
\texttt{topRule} and \texttt{congElem} without the
\texttt{hnotcollection} hypothesis, delegating to the appropriate
bridge for each constructor.

\paragraph{Source-rule provenance boundary.}
The exec rule constructed for \texttt{congElem} is ad-hoc---it is
\emph{not} drawn from \texttt{languageDefToExecRules}.  Source-aware
bridges (\texttt{DeclReducesWithPremises} $\to$ \texttt{SourceExecRule})
retain \texttt{hnotcoll} because their conclusions require a witness in
\texttt{languageDefToSourceExecRules lang}.  This is a genuine
source-language boundary, not a proof gap: encoding \texttt{congElem}
as a \texttt{SourceExecRule} would require a path/focus-aware
source-language extension (Restriction~\ref{sec:future}).

\subsection{Compound Pattern Round-Trips}
\label{sec:matchatom-roundtrip}

The space-effect rules for \texttt{add-atom} and \texttt{remove-atom}
(\texttt{SpaceEffectFragment.lean}, 193~lines) use compound command
patterns like \texttt{(add-atom \&self \$x)}.  Previously, the
correctness of matching these patterns against concrete command
expressions was an MM2 compiler invariant expressed as a hypothesis.

Using the \texttt{MatchAtomRel} relational specification
(\texttt{MatchSpec.lean}), we now construct explicit derivations
(\texttt{expr\_cons} / \texttt{symbol} / \texttt{var\_fresh}) and
apply \texttt{matchAtom\_complete} to obtain proven round-trip
theorems:

\begin{theorem}[\texttt{matchAtom\_addAtomCmdAtom}]
For any \MeTTaIL{} pattern $p$,
\texttt{matchAtom [] addAtomCmdAtom (morkPatternToAtom (.apply "add-atom"
[.apply "\&self" [], p]))} = \texttt{some [("x", morkPatternToAtom p)]}.
\end{theorem}

\noindent
These are then lifted to source-rule firing level:

\begin{theorem}[\texttt{addAtom\_fireSourceRule\_mem}]
When the command atom is in the workspace, \texttt{addAtomSourceExecRule}
fires and produces the \texttt{applySinks} result with binding
\texttt{"x"~$\mapsto$~morkPatternToAtom~p}.
\end{theorem}

\noindent
Symmetric theorems hold for \texttt{remove-atom}.

\subsection{Extended Bridge: Beyond \texttt{morkTranslatable}}
\label{sec:extendedbridge}

The base \MeTTaIL{} bridge requires \texttt{morkTranslatable r.right}
(no \texttt{.subst} nodes, no rest-variable collections).  The extended
bridge (\texttt{ExtendedBridge.lean}, 83~lines, plus supporting
infrastructure in \texttt{MeTTaILBridge.lean} and
\texttt{Substitution.lean}) lifts this restriction by proving that
\texttt{applyBindings} always normalizes these constructs away.

\paragraph{Pattern normalization.}
A mutual recursive function \texttt{normalize} / \texttt{normalizeList}
eliminates all \texttt{.subst} nodes by evaluating them via
\texttt{openBVar}.  The key property:

\begin{theorem}[\texttt{noExplicitSubst\_normalize}]
For all patterns $p$,
\texttt{noExplicitSubst (normalize p) = true}.
\end{theorem}

\paragraph{\texttt{morkTranslatable} preservation under \texttt{openBVar}.}
The theorem \texttt{morkTranslatable\_openBVar} proves that substituting
a translatable term into a translatable pattern preserves translatability.
This is proven by structural induction on \texttt{Pattern.inductionOn},
with the \texttt{.subst} case vacuously true (no translatable
\texttt{.subst} exists) and the \texttt{.collection} case requiring a
match on the rest field.

\paragraph{Core theorem.}
\begin{theorem}[\texttt{morkTranslatable\_applyBindings}]
\label{thm:applyBindings}
For any bindings $\mathit{bs}$ and pattern $\mathit{rhs}$, if
all values in $\mathit{bs}$ are \texttt{morkTranslatable}, then
\texttt{morkTranslatable (applyBindings bs rhs) = true}.
\end{theorem}

\noindent
Note: no \texttt{morkTranslatable rhs} hypothesis is needed.
\texttt{applyBindings} eliminates \texttt{.subst} (via \texttt{openBVar})
and resolves rest-variables (via lookup and splice with
\texttt{rest = none}), so the output is always in the
\texttt{morkTranslatable} fragment regardless of the input structure.

The proof handles each \texttt{Pattern} constructor:
\texttt{.fvar}~$x$ looks up in bindings (translatable by hypothesis
or identity); \texttt{.subst} recurses on body and replacement, then
applies \texttt{morkTranslatable\_openBVar}; \texttt{.collection}
maps elements (IH), extracts rest-variable elements from bindings
(translatable since the bound collection value is translatable),
and appends via \texttt{morkTranslatableList\_append}.

\paragraph{Extended fire theorems.}
\begin{theorem}[\texttt{declReduces\_extended\_mork\_fire}]
A \MeTTaIL{} rewrite step fires via an ad-hoc exec rule even when
the rule's RHS uses \texttt{.subst} or rest-variables.  The witness
uses \texttt{collectionReplaceRule} with groundness of both old and
new atoms.
\end{theorem}

\noindent
A source-rule variant \texttt{declReduces\_extended\_mork\_sourceRuleFire}
provides the same guarantee at the \texttt{SourceExecRule} /
\texttt{fireSourceRule} level.

%% ============================================================
\section{MQ-Calculus Bridge}
\label{sec:mqbridge}
%% ============================================================

The file \texttt{MORKCommBridge.lean} formalizes the correspondence between
\MORK{}'s binary fold steps and the MQ-calculus \texttt{CommReduction} rule.

\begin{theorem}[\texttt{mork\_fold\_is\_comm}]
For each binary fold outcome, there exists an MQ \texttt{CommReduction} with
the matching branch selection.
\end{theorem}

\begin{theorem}[\texttt{mork\_fold\_both\_outcomes\_exist}]
Both outcomes are constructively possible for any binary fold step:
non-determinism is structural.
\end{theorem}

\begin{theorem}[\texttt{mork\_mq\_nondeterminism\_corresponds}]
\MORK{}'s binary fold non-determinism (both sub-results reachable) is
equivalent to MQ's \texttt{comm\_both\_outcomes}.
\end{theorem}

\subsection{N-ary Generalization}

The binary case is a special case of N-ary fold non-determinism:

\begin{theorem}[\texttt{nary\_fold\_all\_outcomes\_exist}]
For any fold step with $N$ sub-results and any index $k < N$, there
exists a fold step selecting sub-result~$k$ as the assembled output.
\end{theorem}

\begin{theorem}[\texttt{binary\_subResult0\_is\_nary}]
Binary \texttt{FoldPicksSubResult} for sub-result~0 is equivalent to
\texttt{NaryFoldPicksSubResult} at index~0.
\end{theorem}

\begin{theorem}[\texttt{nary\_fold\_has\_binary\_comm}]
Any N-ary fold with $\geq 2$ sub-results has both COMM outcomes available.
\end{theorem}

\noindent
Seven canary theorems protect these invariants.

%% ============================================================
\section{Conformance Testing}
\label{sec:conformance}
%% ============================================================

The conformance test suite (\texttt{Conformance.lean}, 1{,}407 lines)
provides 13 kernel-checked conformance theorems verified against the
\MORK{} CLI (\texttt{mork run}), plus bidirectional match-pattern and
fire-rule correspondence theorems between the computable and spec levels.

\subsection{Computable Reference Evaluator}

The spec-level \texttt{fireRule} is \texttt{noncomputable} (uses
\texttt{Finset.toList}).  The computable reference evaluator mirrors
these operations over \texttt{List Atom}:

\begin{itemize}[nosep]
  \item \texttt{cmatchAtom} / \texttt{cmatchAtomList}: mutual recursive
    matcher that extends \texttt{matchAtom} to handle \texttt{.expression}
    patterns by element-wise matching.
  \item \texttt{cmatchPattern}: conjunctive matching over a list space.
  \item \texttt{capplySinks}: sink application on list spaces (handles
    \texttt{add}, \texttt{remove}, and \texttt{head} with explicit
    membership check for idempotence).
  \item \texttt{cfireRule}: end-to-end rule firing.
\end{itemize}

\noindent
All operations reduce by \texttt{rfl} in the Lean kernel.

\subsection{Test Fixtures}

Each test is defined as an \texttt{ExecRule} applied to a concrete space,
with a theorem asserting that \texttt{cfireRule} produces the expected
output.  The corresponding \texttt{.mm2} file is run through
\texttt{mork run} to obtain the ground-truth output.

\begin{table}[ht]
\centering
\caption{Conformance test suite.  All theorems are proved by \texttt{rfl}.}
\label{tab:conformance}
\begin{tabular}{@{}clll@{}}
\toprule
\textbf{\#} & \textbf{Feature} & \textbf{MM2 Source} & \textbf{Theorem} \\
\midrule
1 & Expression add/remove  & \texttt{test\_add\_simple.mm2} & \texttt{conformance\_test1\_add\_simple} \\
2 & Flat symbol add        & \texttt{test\_add\_constant.mm2} & \texttt{conformance\_test2\_add\_constant} \\
3 & Variable binding       & \texttt{test3\_var\_binding.mm2} & \texttt{conformance\_test3\_var\_binding} \\
4 & Conjunctive (shared var) & \texttt{test4\_conjunctive.mm2} & \texttt{conformance\_test4\_conjunctive} \\
5 & Equality (\texttt{\$x \$x})  & \texttt{test5\_equal\_pair.mm2} & \texttt{conformance\_test5\_equal\_pair} \\
6 & Pattern mismatch       & \texttt{test6\_no\_match.mm2} & \texttt{conformance\_test6\_no\_match} \\
7 & Nested expressions     & \texttt{test7\_nested.mm2} & \texttt{conformance\_test7\_nested} \\
8 & Multi-step fixpoint    & \texttt{test8\_multi\_step.mm2} & \texttt{conformance\_test8\_*} (5 thms) \\
9 & Arity mismatch         & (inline) & \texttt{conformance\_test9\_arity\_mismatch} \\
\bottomrule
\end{tabular}
\end{table}

\noindent
Test~8 exercises multi-step exhaustive execution: the rule fires three times,
consuming each \texttt{(task $x$)} atom and producing \texttt{(done $x$)}.
The formalization proves each step individually via \texttt{rfl} (the
fixpoint function does not reduce definitionally, so a chain of one-step
equalities is used instead).

\subsection{Expression Matching}

The spec-level \texttt{matchAtom} returns \texttt{none} on
\texttt{.expression} patterns.  The computable evaluator's
\texttt{cmatchAtom} extends this via mutual recursion:

\begin{lstlisting}[caption={Expression matching (\texttt{Conformance.lean}).}]
mutual
def cmatchAtom (sig : Subst) (pat conc : Atom)
    : Option Subst :=
  match pat, conc with
  | .expression ps, .expression cs =>
    cmatchAtomList sig ps cs
  | .var v, a => ...  -- same as matchAtom
  | ...
def cmatchAtomList (sig : Subst) (pats concs
    : List Atom) : Option Subst :=
  match pats, concs with
  | [], [] => some sig
  | p :: ps, c :: cs =>
    match cmatchAtom sig p c with
    | some sig' => cmatchAtomList sig' ps cs
    | none => none
  | _, _ => none
end
\end{lstlisting}

\noindent
This enables conformance tests with real MM2 S-expressions (tests 1, 3--8).

%% ============================================================
\section{Coverage and Open Gaps}
\label{sec:gaps}
%% ============================================================

Table~\ref{tab:coverage} summarizes the boundary between formalized and
unformalized aspects.

\begin{table}[ht]
\centering
\caption{Formal coverage boundary.}
\label{tab:coverage}
\begin{tabular}{@{}lll@{}}
\toprule
\textbf{Aspect} & \textbf{Formalized} & \textbf{Not Formalized} \\
\midrule
MM2 syntax                     & add/remove/head sinks, FoldAggregator, SourceFactor & Remaining 11 Rust sink types \\
Pattern matching (spec)        & Soundness + completeness (expressions) & Coreferential matching \\
Pattern matching (evaluator)   & Expressions, nested, variables & --- \\
Space transitions              & Yes + sink persistence & Nested collection rewrite \\
Three-phase protocol           & Priority bands + aggregator & --- \\
Work-queue scheduler           & Read-copy, self-respawn, ordering, full correspondence & Concurrent scheduling \\
Source-side input              & BTM, == (equality), != (inequality) & ACT (external file), z3 \\
Conformance testing            & 24 kernel-checked theorems, 9+ MM2 tests & --- \\
\PathMap{} algebra             & Quantale hierarchy + laws & --- \\
\PathMap{} trie (\texttt{FTrie}) & Join/meet/subtract/restrict + refinement & Full commutativity (non-Unit) \\
\MORK{}--\PathMap{} bridge     & Algebraic (join/subtract) & --- \\
\MeTTaIL{} bridge              & Flat rules (\texttt{.fvar}), multi-premise, freshness guards, \texttt{congElem} (zipper/lens + ad-hoc), \texttt{LensRel} laws, \texttt{matchAtom} round-trips for compound patterns, extended bridge (\texttt{.subst}/rest-var resolution via \texttt{applyBindings}) & Source-level positional congruence \\
PLN routing witnesses          & 6 rules (transitivity, MP, abduction, guarded) & --- \\
MQ bridge                      & Binary + N-ary fold $\leftrightarrow$ COMM & Real process witnesses \\
Execution boundary             & \texttt{morkTranslatable} packaging, pattern normalization (\texttt{normalize}), \texttt{morkTranslatable\_applyBindings} & --- \\
Runtime kernel packaging       & Effect-class bridge, classified fragments (exec/query/spaceEffect), memoization safety, resource descriptors (4 resources) & Named atomspace exec, Z3/ACT exec seams \\
\bottomrule
\end{tabular}
\end{table}

\subsection{Notable Restrictions}

\begin{enumerate}
  \item \textbf{Expression matching.}
    The spec-level \texttt{matchAtom} now handles expression patterns
    (via mutual recursion with \texttt{matchAtomList}).
    The biconditional \texttt{matchAtom\_iff} covers the
    \texttt{morkTranslatable} fragment (excluding \texttt{.subst} and
    rest-variable \texttt{.collection}).  However,
    \texttt{morkTranslatable\_applyBindings}
    (\S\ref{sec:extendedbridge}) proves that \texttt{applyBindings}
    eliminates both constructs, so rule RHS patterns containing
    \texttt{.subst} or rest-variables are fully handled at the
    bridge level.

  \item \textbf{Coreferential matching.}
    \MORK{}'s Rust backend uses de~Bruijn-style coreferential matching with
    byte capture (\texttt{\$ps}, \texttt{\$cs}).
    The formalization uses first-order variable binding.

  \item \textbf{Aggregation sinks.}
    The formalization captures \texttt{add}, \texttt{remove}, \texttt{head},
    and fold-level \texttt{count}/\texttt{sum}/\texttt{selectFirst}/\texttt{selectAll}.
    The Rust backend's remaining 11 per-match sink types
    (\texttt{min}, \texttt{max}, etc.) are not formalized.

  \item \textbf{Aggregator enforcement.}
    The \texttt{AggregatorConsistent} predicate
    (re-exported via \texttt{ExecutionBoundary.aggregatorConsistency})
    now connects fold-assembled results to aggregator semantics.
    Its use as a runtime invariant is not yet enforced.

  \item \textbf{Computability.}
    The spec-level space operations are \texttt{noncomputable}
    (\texttt{Finset.toList}).  The computable reference evaluator provides
    a kernel-reducible alternative for conformance testing.
    The scheduler correspondence chain
    (\S\ref{sec:schedulercorr}) bridges the gap: under
    \texttt{WorkQueueInvariant}, the computable scheduler agrees with
    the spec through arbitrary bounded-run horizons.

  \item \textbf{FTrie commutativity.}
    \texttt{FTrie} join is left-biased (\texttt{v$_1$ <|> v$_2$ = v$_1$}),
    so \texttt{JoinComm} does not hold in general.
    For \texttt{V = Unit} (\MORK{}'s use case), commutativity is trivial.

  \item \textbf{MQ bridge witnesses.}
    The MQ-MORK bridge uses \texttt{Process.MQNil} witnesses---the
    correspondence is structural but the witnesses carry no information
    about the actual \MORK{} processes.

  \item \textbf{Metatheory.}
    Termination for ground removing programs and confluence for
    deterministic rule sets are not yet formalized.

  \item \textbf{Source-level positional congruence.}
    \label{sec:future}
    The zipper/lens layer provides the structural semantics for
    sub-expression access.  The remaining gap is encoding
    \texttt{congElem} as a genuine \texttt{SourceExecRule} with
    path-aware source factors (\texttt{SourceFactor.focus}) and sink
    actions (\texttt{SinkAction.replaceAtPath}), which would discharge
    the \texttt{hnotcoll} hypothesis from source-aware bridges.
\end{enumerate}

%% ============================================================
\section{Execution Boundary}
\label{sec:execboundary}
%% ============================================================

\texttt{ExecutionBoundary.lean} (529~lines) provides a thin re-export
surface packaging the proven \texttt{morkTranslatable} execution boundary,
the scheduler correspondence chain, the extended bridge infrastructure,
the zipper/lens congruence infrastructure (\texttt{AtomZipper},
\texttt{LensRel}, lens laws), the collection congruence bridge, and
the extended \texttt{.subst}/rest-variable bridge.
The predicate \texttt{morkTranslatable} identifies the \MeTTaIL{} pattern
fragment for which substitution commutation holds.  However, the
extended bridge (\S\ref{sec:extendedbridge}) proves that patterns
\emph{outside} this fragment (containing \texttt{.subst} or rest-variable
collections) are still handled: \texttt{applyBindings} normalizes them
into the fragment, and ad-hoc replace rules provide execution witnesses.

This file re-exports key definitions (\texttt{morkTranslatable},
\texttt{morkPatternToAtom}, pattern-to-space translation, rewrite-rule
translation), scheduler correspondence theorems
(\texttt{cFireExecFact\_toFinset}, \texttt{cWorkQueueStep\_toFinset},
\texttt{cWorkQueueRunN\_toFinset}, \texttt{WorkQueueInvariant}),
source guard infrastructure (\texttt{SourceGuard}, \texttt{matchSourceGuard},
\texttt{matchSourceGuards}, \texttt{isAtomFresh},
\texttt{freshness\_premise\_correspond}),
extended bridge theorems (\texttt{premiseChain\-MatchFactorsExt},
\texttt{multiPremiseExtBridge}, \texttt{fullExtBridge},
\texttt{languageDefSourceRulesExt}),
zipper/lens infrastructure (\texttt{AtomZipper}, \texttt{AtomCrumb},
\texttt{LensRel}, lens laws, \texttt{collection\_lensRel}),
collection congruence bridge theorems
(\texttt{congElemZipperBridge}, \texttt{declReducesFullBridge},
\texttt{groundSelfMatch}, \texttt{collectionReplaceFires},
\texttt{congElemBridge}),
and the extended bridge (\texttt{patternNormalize},
\texttt{normalizeSubstFree},
\texttt{openBVarPreservesTranslatable},
\texttt{applyBindingsPreservesTranslatable},
\texttt{extendedExecFire},
\texttt{extendedSourceRuleFire})
as public \texttt{abbrev}s, providing a stable API surface
for downstream consumers.

%% ============================================================
\section{Runtime-Kernel Packaging}
\label{sec:runtimekernel}
%% ============================================================

Two companion files above the \MORK{} layer package the current proved
execution seams into a backend-neutral runtime-kernel object that the
\texttt{ElaboratedCore} elaboration layer can target without importing
scheduler internals.

\subsection{Effect-Class Bridge (\texttt{RuntimeKernel.lean})}

\texttt{RuntimeKernel.lean} (211~lines) bridges the \texttt{RuntimeKernelClass}
enum (from \texttt{ElaboratedCoreBase.lean}) to the \texttt{EffectClass}
hierarchy (from \texttt{mettail-core/EffectSafety.lean}).

\begin{definition}[\texttt{RuntimeKernelClass.effectClass}]
\begin{tabular}{@{}ll@{}}
\texttt{.ruleExec} & $\mapsto$ \texttt{pureStructural} \\
\texttt{.query} & $\mapsto$ \texttt{readOnlyLookup} \\
\texttt{.spaceEffect} & $\mapsto$ \texttt{writesState} \\
\texttt{.oracle} & $\mapsto$ \texttt{oracleIO} \\
\texttt{.metaPhase} & $\mapsto$ \texttt{pureStructural}
\end{tabular}
\end{definition}

\noindent
A \texttt{ClassifiedFragment} pairs a kernel class with its effect class,
resource class, and backend name.
Three canonical fragments package the current \MORK{}/MM2-backed witnesses:
\begin{itemize}[nosep]
  \item \texttt{execFragment}: eval/evalc $\to$ \texttt{pureStructural} on \texttt{defaultAtomSpace}
  \item \texttt{queryFragment}: match/get-atoms $\to$ \texttt{readOnlyLookup} on \texttt{defaultAtomSpace}
  \item \texttt{spaceEffectFragment}: add-atom/remove-atom $\to$ \texttt{writesState} on \texttt{defaultAtomSpace}
\end{itemize}

\noindent
A \texttt{RuntimeKernelPackage} bundles the \texttt{MeTTaRuntimeKernelTriad}
with these classified fragments.  The canonical instance
\texttt{morkRuntimeKernelPackage} uses all three.

Memoization safety is kernel-checked:
exec and query fragments support both scalar and outcome-set memoization;
\texttt{spaceEffect} supports neither (by \texttt{decide}).
All theorems have zero axioms.

\subsection{Resource Layer (\texttt{RuntimeResource.lean})}

\texttt{RuntimeResource.lean} (234~lines) adds an explicit resource layer
classifying which resources the current \MORK{}/MM2 backend supports.

\begin{definition}[\texttt{ResourceDescriptor}]
A descriptor carries: resource class, supported kernel classes, effect
envelope, backend name, and a Boolean \texttt{hasProvedExecSeam}.
\end{definition}

\noindent
Four descriptors are defined, aligned with the \MORK{} Rust implementation:

\begin{table}[ht]
\centering
\caption{Runtime resource descriptors.}
\label{tab:resources}
\begin{tabular}{@{}lllll@{}}
\toprule
\textbf{Resource} & \textbf{Kernel Classes} & \textbf{Effect} & \textbf{Backend} & \textbf{Proved?} \\
\midrule
\texttt{defaultAtomSpace}  & ruleExec, query, spaceEffect & writesState & MORK/MM2 & yes \\
\texttt{namedAtomSpace}    & (none)                       & writesState & none     & no  \\
\texttt{solverResource}    & oracle                       & oracleIO    & MORK/MM2 & no  \\
\texttt{externalResource}  & oracle                       & oracleIO    & MORK/MM2 & no  \\
\bottomrule
\end{tabular}
\end{table}

\noindent
The \texttt{solverResource} and \texttt{externalResource} descriptors
acknowledge that \MORK{}'s Rust runtime actively uses Z3 subprocesses
(\texttt{ResourceRequest::Z3} in \texttt{space.rs}) and ACT
file-backed external datasets, but their Lean formalization remains
future work (\texttt{Syntax.lean}: ``NOT yet formalized: ACT, z3'').
Named atomspaces are not supported by \MORK{}'s single-space-per-instance
architecture.

A \texttt{ResourceProfile} pairs a \texttt{RuntimeKernelPackage} with its
resource descriptors and identifies the primary resource.
The canonical \texttt{morkResourceProfile} uses \texttt{defaultAtomSpace}
as the primary resource (the only one with a proved execution seam).

All classification theorems (19 in \texttt{RuntimeKernel.lean},
16 in \texttt{RuntimeResource.lean}) are proved by \texttt{rfl} or
\texttt{decide} with axioms at most \texttt{[propext, Classical.choice,
Quot.sound]}.

%% ============================================================
\section*{Acknowledgements}
%% ============================================================

The \MORK{}/MM2 design is due to Meredith, Stay, and
collaborators~\cite{MeredithStayOSLF2020}.
Luke Peterson designed and implemented the \MORK{} kernel and MM2 language.
Adam Vandervorst created the \PathMap{} data structure and its Rust
implementation.
Remy Clarke contributed MM2 structuring patterns and example programs.
The Lean~4 formalization was developed with AI assistance (Claude Code,
ChatGPT, Codex).

\begin{thebibliography}{9}

\bibitem{MeredithStayOSLF2020}
L.~G.~Meredith, M.~Stay, et~al.
\newblock \emph{MORK: \MeTTa{} Optimal Reduction Kernel}.
\newblock Design notes and Rust implementation, 2020--2026.

\bibitem{ProcessCalculiPaper}
Zar and Oru\v{z}i.
\newblock Process calculi formalized in {Lean}~4: {Rho}, {Pi}, {MQ}, and
  modal mu-calculus with cross-calculus bridges.
\newblock Technical note, March 2026.

\end{thebibliography}

\end{document}
